% problem-000102 8 有机反应机理
\section*{8. 有机反应机理}
\textit{下列为分问。}

\subsection*{8-1}
如下光学纯(�)-对溴苯磺酸酯在75%的1,4-⼆氧六环⽔溶液中与 $\mathrm { N a N _ { 3 } }$ 反应，得到构型保持产物(�)-2-叠氮⾟烷。画出该反应中间体的结构简式。

（图：）

\subsection*{8-2}
如下光学纯桥环化合物1在酸性条件下与⼄酸酐反应，得到外消旋的酯2。画出该反应关键中间体的结构简式。

（图：）

\subsection*{8-3}
化合物3是⼀种“受阻Lewis酸碱对”，能够在室温下活化氢⽓并形成4；当加热到 $1 5 0 ~ ^ { \circ } \mathrm { C }$ 时，4可释放出氢⽓。4与苯甲醛在⼆氯甲烷溶液中反应⽣成5。画出4和5的结构结构简式。

（图：）

\subsection*{8-4}
螺环化合物6与苯甲醛在邻⼆甲苯中加热回流，反应⽣成化合物7和8。画出产物7以及反应中间体的结构简式。

（图：）

\subsection*{8-5}
霉酚酸的全合成关键中间体13的合成路线如下。画出中间产物9、10、11和12的结构简式。

（图：）

$
\left(\mathrm{TBDMS} = t - \mathrm{BuMe} _ {2} \mathrm{Si}; n - \mathrm{BuLi} = n - \mathrm{C} _ {4} \mathrm{H} _ {9} \mathrm{Li}; \mathrm{DMAP} = 4 - \text {二甲氨基吡啶}; \mathrm{Ms} = \mathrm{MeSO} _ {2}\right)
$
